\chapter{Application Shells}

Once you can compose widgets, the next layer is the application shell: the main window, menu bar, toolbars, status bar, dock widgets, and reusable commands. In Qt, this shell is centered on \api{QMainWindow}. In \api{crystal-qt6}, the wrapper is \api{Qt6::MainWindow}.

\section{Main Window Basics}

\api{MainWindow} is a specialized top-level window. It owns a central widget and reserves standard regions for menus, toolbars, docks, and a status bar.

\begin{lstlisting}[style=crystal]
app = Qt6.application

main = Qt6::MainWindow.new
main.window_title = "Map Tool"
main.resize(960, 640)

canvas = Qt6::Label.new("Canvas")
main.central_widget = canvas

main.status_bar.show_message("Ready")
main.show
app.run
\end{lstlisting}

The central widget is the primary work area. In an editor, it might be a custom canvas. In a document tool, it might be a text editor. In a data tool, it might be a table or a split panel.

\section{Actions First}

Qt applications usually model user commands as actions. An action can have text, a shortcut, enabled state, checked state, tooltip text, status-tip text, and callbacks. The same action can appear in a menu and a toolbar without duplicating behavior.

\begin{lstlisting}[style=crystal]
open_action = Qt6::Action.new("Open", main)
open_action.shortcut = "Ctrl+O"
open_action.tool_tip = "Open a project"
open_action.status_tip = "Open a project file"
open_action.on_triggered do
  main.status_bar.show_message("Open requested", 1500)
end
\end{lstlisting}

This pattern scales better than wiring every menu item and toolbar button separately. Menus and toolbars become views of the same command model.

\section{Shortcuts}

Shortcuts can be assigned as strings or explicit \api{KeySequence} values. The string form is concise and works well for common shortcuts:

\begin{lstlisting}[style=crystal]
save_action = Qt6::Action.new("Save", main)
save_action.shortcut = "Ctrl+S"

apply_action = Qt6::Action.new("Apply", main)
apply_action.shortcut = Qt6::KeySequence.new("Ctrl+Return")
\end{lstlisting}

Prefer familiar platform conventions. Use shortcuts for commands that users repeat often: open, save, undo, redo, find, apply, close, and mode switching.

\section{Menus}

The main window owns a menu bar. Top-level menus are created from the menu bar, and each menu can contain actions, separators, and submenus.

\begin{lstlisting}[style=crystal]
file_menu = main.menu_bar.add_menu("File")
file_menu << open_action
file_menu << save_action
file_menu.add_separator

quit_action = Qt6::Action.new("Quit", main)
quit_action.shortcut = "Ctrl+Q"
quit_action.on_triggered { app.quit }
file_menu << quit_action

view_menu = main.menu_bar.add_menu("View")
\end{lstlisting}

Menus should be organized by intent. \api{File} usually holds document or project commands. \api{Edit} holds undo, redo, cut, copy, paste, and selection commands. \api{View} controls panes and presentation. \api{Tools} or \api{Mode} can hold domain-specific commands.

\section{Toolbars}

Toolbars expose frequent commands as visible controls. In \api{crystal-qt6}, a toolbar can contain actions, separators, and ordinary widgets.

\begin{lstlisting}[style=crystal]
toolbar = Qt6::ToolBar.new("Main", main)
toolbar << open_action
toolbar << save_action
toolbar.add_separator
toolbar << apply_action
toolbar.movable = true

main.add_tool_bar(toolbar)
\end{lstlisting}

Toolbar actions should be a curated subset of menu actions. A toolbar is most useful when it keeps high-frequency operations close to the work area without replacing the menu bar as the complete command catalog.

\section{Status Bars}

The status bar is for quiet, transient feedback. Use it for command results, validation summaries, progress hints, selection details, and mode changes.

\begin{lstlisting}[style=crystal]
status = main.status_bar
status.show_message("Ready")

save_action.on_triggered do
  status.show_message("Saved project", 2000)
end
\end{lstlisting}

A timeout of \api{0} keeps the message until it is replaced or cleared. Nonzero timeouts are useful for short command confirmations.

\section{Docks}

Dock widgets attach secondary panels to a main window. They are a natural fit for inspectors, layer lists, logs, search results, palettes, and tool options.

\begin{lstlisting}[style=crystal]
layers = Qt6::ListWidget.new
layers << "Terrain" << "Roads" << "Units"

layers_dock = Qt6::DockWidget.new("Layers", main)
layers_dock.widget = layers
main.add_dock_widget(layers_dock, Qt6::DockArea::Left)

properties = Qt6::Widget.new
properties.form do |form|
  form.add_row("Name", Qt6::LineEdit.new("Terrain"))
  form.add_row("Visible", Qt6::CheckBox.new("Visible"))
end

properties_dock = Qt6::DockWidget.new("Properties", main)
properties_dock.widget = properties
main.add_dock_widget(properties_dock, Qt6::DockArea::Right)
\end{lstlisting}

Keep docks focused. A dock should have one clear purpose, and its title should make that purpose obvious.

\section{Checkable Actions And Mode Groups}

Actions can be checkable. Related checkable actions can be placed in an \api{ActionGroup}. When the group is exclusive, only one action can be checked at a time.

\begin{lstlisting}[style=crystal]
mode_group = Qt6::ActionGroup.new(main)
mode_group.exclusive = true

select_mode = Qt6::Action.new("Select", main)
select_mode.checkable = true
select_mode.checked = true

draw_mode = Qt6::Action.new("Draw", main)
draw_mode.checkable = true

pan_mode = Qt6::Action.new("Pan", main)
pan_mode.checkable = true

mode_group << select_mode
mode_group << draw_mode
mode_group << pan_mode
\end{lstlisting}

Mode actions can appear in both menus and toolbars:

\begin{lstlisting}[style=crystal]
mode_menu = main.menu_bar.add_menu("Mode")
mode_menu << select_mode
mode_menu << draw_mode
mode_menu << pan_mode

toolbar.add_separator
toolbar << select_mode
toolbar << draw_mode
toolbar << pan_mode
\end{lstlisting}

\section{A Small Shell}

The following shell combines the main pieces: central widget, actions, shortcuts, menus, toolbar, status bar, and a dock.

\begin{lstlisting}[style=crystal]
require "qt6"

app = Qt6.application

main = Qt6::MainWindow.new
main.window_title = "Project Shell"
main.resize(900, 600)

status = main.status_bar
status.show_message("Ready")

canvas = Qt6::Label.new("Project canvas")
canvas.set_minimum_size(480, 320)
main.central_widget = canvas

open_action = Qt6::Action.new("Open", main)
open_action.shortcut = "Ctrl+O"
open_action.status_tip = "Open a project"
open_action.on_triggered do
  status.show_message("Open requested", 1500)
end

save_action = Qt6::Action.new("Save", main)
save_action.shortcut = "Ctrl+S"
save_action.status_tip = "Save the current project"
save_action.on_triggered do
  status.show_message("Saved", 1500)
end

quit_action = Qt6::Action.new("Quit", main)
quit_action.shortcut = "Ctrl+Q"
quit_action.on_triggered { app.quit }

file_menu = main.menu_bar.add_menu("File")
file_menu << open_action
file_menu << save_action
file_menu.add_separator
file_menu << quit_action

toolbar = Qt6::ToolBar.new("Main", main)
toolbar << open_action
toolbar << save_action
main.add_tool_bar(toolbar)

inspector = Qt6::Widget.new
inspector.form do |form|
  form.add_row("Name", Qt6::LineEdit.new("Terrain"))
  form.add_row("Visible", Qt6::CheckBox.new("Visible"))
end

dock = Qt6::DockWidget.new("Inspector", main)
dock.widget = inspector
main.add_dock_widget(dock, Qt6::DockArea::Right)

main.show
app.run
\end{lstlisting}

\section{Shell Design Guidelines}

A good shell gives users stable places to look:

\begin{enumerate}
\item Put the primary work surface in the central widget.
\item Put complete command access in menus.
\item Put frequent commands in toolbars.
\item Put persistent secondary tools in docks.
\item Put transient feedback in the status bar.
\item Reuse actions instead of duplicating callbacks.
\end{enumerate}

That structure is familiar to users of desktop applications, and it keeps the code organized as the application grows.

\section{Examples}

\api{examples/editor_shell.cr} is the smallest maintained example focused on the main-window shell. \api{examples/desktop_editor_showcase.cr} is broader and shows how shell structure combines with model/view panels, text editing, clipboard data, image loading, custom preview rendering, and application-service helpers.
